\documentclass[11pt,a4paper]{article}
\usepackage{../pelm}
\begin{document}
\paper{Quiz 1}{Weeks 1--2: Introduction, Vocabulary, and the History of AI}{10}{10 minutes}

\begin{questions}

\q{Which of the following is an example of \textbf{generative} AI?}
  {A bank system deciding whether a transaction is suspicious}
  {A hospital system deciding which scans a radiologist should read first}
  {A system that drafts a letter to a patient's family}
  {A spam filter sorting incoming email}

\q{In this course, the term \textbf{prompt} refers to:}
  {Only the question you type}
  {Everything you give the system, including attached documents and images}
  {The system's internal instructions, which you cannot see}
  {The first message of a conversation only}

\q{What is the \textbf{context}?}
  {Everything the model learned during training}
  {Everything the model can see at the moment it answers you}
  {The subject area a question belongs to}
  {The model's permanent memory of your previous conversations}

\q{A student asks an AI assistant about employment notice periods in Turkey and receives a
   confident, well-structured answer that actually describes United States practice. Which
   failure signature is this?}
  {Fabrication}
  {Displacement}
  {Omission}
  {Overconfidence}

\q{An AI-generated summary contains a citation in perfect format, attributed to a real journal,
   for an article that does not exist. Which failure signature is this?}
  {Fabrication}
  {Displacement}
  {Omission}
  {Silent failure}

\q{Which failure signature is hardest to detect, and why?}
  {Fabrication, because citations are difficult to check}
  {Displacement, because foreign law is complex}
  {Omission, because there is no wrong sentence to point at}
  {Overconfidence, because the tone is always the same}

\q{ELIZA (1966) is significant to this course mainly because:}
  {It was the first program to pass the Turing test}
  {It demonstrated that fluent language alone is enough to produce a sense of being understood}
  {It was the first system to use neural networks}
  {It proved that rule-based approaches could scale}

\q{The ``ELIZA effect'' describes:}
  {The tendency of chatbots to repeat themselves}
  {The tendency of people to attribute understanding to a system producing plausible language}
  {The loss of accuracy when a model is made smaller}
  {The tendency of models to agree with the user}

\q{Herbert Simon predicted in 1965 that machines would be able to do any work a human can do
   within twenty years. What lesson does this course draw from such predictions?}
  {Predictions about AI are always wrong}
  {The direction of predictions has often been reasonable while the timing has been badly optimistic}
  {AI progress has been faster than predicted}
  {Researchers should not make predictions}

\q{Across the four eras of the field (rules, statistics, neural networks, transformers), what is
   the constant direction of change?}
  {From smaller models towards larger models}
  {From English towards multilingual systems}
  {From human-written specification towards learning from data}
  {From academic research towards commercial products}

\end{questions}
\end{document}
